\chapter{Detection, Intervention, and Pre-Boarding Screening}
\section{Scientific claim}
Mid-voyage response saturates early, while pre-boarding screening has a larger decision value
because it acts before infectious passengers or crew enter the shipboard contact network.
In C14 and C14b, syndromic or cascade response reduced final attack rates for all modeled
pathogens, and wearables, wastewater, and multiplex testing added less than one percentage point
for most controlled comparisons.
The boundary decision model changes the question from ``which signal detects an outbreak first''
to ``which pre-boarding policy reduces the expected cost of a voyage.''
On that criterion, incentivized wearable adoption followed by confirmatory testing is the preferred
active policy across most prevalence, pathogen, and platform cells.

\section{Mid-voyage surveillance results: response saturation}
The active-outbreak campaigns compared \nonetrue{} against syndromic response, cascade response,
cascade plus multiplex testing, wearable variants, and wastewater surveillance.
Table~\ref{tab:ch7-intervention-comparison} consolidates the cross-pathogen result.
The uncontrolled attack rate is the mean attack rate under \nonetrue{}.
The controlled range reports the span across syndromic and advanced response strategies in the
C14/C14b intervention summaries.

\begin{table}[htbp]
\centering
\caption{Cross-pathogen intervention comparison for active mid-voyage response. Attack rates are percentages. The advanced marginal gain is the best advanced strategy relative to syndromic response; negative values indicate that syndromic response was already the lowest mean in the supplied summary.}
\label{tab:ch7-intervention-comparison}
\scriptsize
\begin{tabular}{p{1.8cm}rrrrp{2.8cm}}
\toprule
Pathogen & Uncontrolled AR & Syndromic AR & Controlled range & Best advanced gain & Interpretation \\
\midrule
Norovirus & 17.21 & 5.87 & 5.23--5.87 & 0.64 & Advanced signals add $<1$ pp \\
SARS-CoV-2 & 20.72 & 0.38 & 0.38--0.86 & 0.00 & Response almost eliminates onboard spread \\
Measles & 37.91 & 1.32 & 1.15--2.17 & 0.17 & High immunity leaves little preventable residual \\
Influenza & 66.08 & 1.89 & 1.20--2.14 & 0.69 & Fast control after first trigger \\
\emph{C. difficile} & 12.03 & 0.18 & 0.18--0.55 & 0.00 & Clinical response dominates tested add-ons \\
\bottomrule
\end{tabular}
\end{table}

Norovirus is the limiting enteric case.
The uncontrolled mean was 17.21\%, syndromic response reduced it to 5.87\%, and the best advanced
mean in the supplied comparison was 5.23\%.
That 0.64 percentage-point difference is small compared with the 11.34 percentage-point reduction
obtained by activating response at all.
For SARS-CoV-2, the analogous contrast is 20.72\% to 0.38\% under syndromic response, with advanced
strategies spanning 0.49--0.86\% in this summary.
Influenza begins at a much higher uncontrolled attack rate, 66.08\%, but falls to 1.20--2.14\%
after response.

This saturation result should be interpreted operationally rather than as a negative result for
wearable or wastewater instruments.
Once the voyage has entered an active response state, intensified isolation, sanitation, and case
management remove much of the remaining preventable transmission.
Additional signals then compete for a small residual benefit.
That pattern is consistent with the cruise-ship norovirus literature, where the timing of reporting
and response is a major determinant of outbreak size \citep{Jenkins2021,Mouchtouri2024,CDC_VSP2025}.
It is also consistent with respiratory cruise outbreaks, where late detection after shipboard
amplification can leave few low-cost options \citep{Yamagishi2020,Willebrand2022}.

\section{Why boundary surveillance matters more than mid-voyage surveillance}
Boundary surveillance acts on introductions rather than on transmission that has already occurred.
In a mid-voyage setting, a wearable flag, wastewater signal, or multiplex result can only change the
course of an outbreak after at least one infectious person has boarded, generated contacts, and
possibly triggered secondary cases.
At pre-boarding, the same signal can reduce the number of infectious people who board.
The relevant state variable is therefore $K_{\mathrm{board}}$, the number of infectious passengers
and crew after screening, rather than the time of the first onboard alert.

The economic objective also changes at the boundary.
For a cruise operator, the value of screening is not only the number of cases prevented.
It includes avoidance of Vessel Sanitation Program escalation, voyage disruption, regulatory
reporting, cleaning surges, denied-boarding compensation, and reputational loss.
The decision model scores a policy by expected voyage cost, then reports the value of information per passenger as
\begin{equation}
\mathrm{VoI}_{\mathrm{pax}} = \frac{E[C_{P0}] - E[C_{p}]}{N_{\mathrm{pax}}},
\end{equation}
where $P0$ is no screening, $p$ is the screening policy, and $N_{\mathrm{pax}}$ is the passenger
capacity for the platform.
Positive values mean the policy reduces expected cost relative to no screening.

\section{Pre-boarding decision model}
The Monte Carlo decision model combines four inputs: pre-boarding prevalence, wearable adoption and
accuracy, confirmatory testing, and a posterior outbreak surface indexed by the number of infectious
individuals who board.
Each scenario draws infectious status for passengers and crew, applies the screening policy, looks
up the corresponding outbreak outcome at $K_{\mathrm{board}}$, and computes expected cost.
The corrected analysis contains 528 scenarios per pathogen: four platform classes, twenty-two
prevalence grid entries, and six policies.
Each scenario uses 2,000 Monte Carlo voyages.

The six policies are:
\begin{description}
\item[P0] No pre-boarding screening.
\item[P1] Advisory-only screening. In the implemented model, this has no operational screening
charge and does not change boarding, so it is a same-path baseline for P0.
\item[P2] Bring-your-own-device wearable adoption with confirmatory testing; boarding is denied only
after a positive confirmatory test.
\item[P3] Wearable flag alone; boarding is denied after an anomaly without confirmatory testing.
\item[P4] Incentivized wearable adoption, with at least 70\% passenger adoption, followed by
confirmatory testing.
\item[P5] Crew-mandatory adoption with the same confirmatory-test boarding rule as P2.
\end{description}

Wearable sensitivity and specificity are pathogen-specific.
The decision runs used sensitivity/specificity of 0.55/0.85 for norovirus, 0.65/0.85 for influenza,
0.70/0.85 for SARS-CoV-2, and 0.75/0.88 for measles.
Confirmatory testing used sensitivity 0.90 and specificity 0.98.
The mid-cost scenario assigned \$1 per wearable screen, \$75 per confirmatory test, \$2,000 per
false-positive or true-positive denied boarding, \$400 per onboard case, \$500,000 per outbreak,
\$2 million for VSP escalation, and \$2 million for reputational cost.
These values are deliberately separable parameters rather than fixed biological assumptions.

\section{Outbreak surfaces from the $k$-sweep campaign}
The boundary model uses $k$-indexed outbreak surfaces generated from 17,600 simulation runs.
Each surface is keyed by platform class, pathogen, baseline response, and the number of infectious
individuals who board.
For each $k$, the surface stores the probability of VSP triggering, expected attack rate, probability
of acceleration, expected onboard cost, expected peak epoch, and run count.
The decision model interpolates these curves when the Monte Carlo draw produces a non-integer lookup
point, although the implemented screening pipeline returns integer $K_{\mathrm{board}}$.

The surfaces explain why screening value rolls over.
When $K_{\mathrm{board}}$ is low, intercepting one infectious individual can move a voyage from a
high-risk point on the surface to a low-risk point.
When prevalence is high, the no-screening voyage is already near the high-risk plateau, and screening
cannot intercept enough infectious people to move the operating point out of saturation.
False positives and screening costs then continue to rise while marginal VSP avoidance falls.
This produces the large-ship rollover in value of information.

Surface shape differs by pathogen.
Influenza reaches VSP-trigger probabilities above 0.5 at $k=1$ and above 0.9 at $k=2$ on all four
platforms in the supplied surface.
Norovirus reaches 0.5 at $k=2$ and 0.9 at $k=4$ or $k=5$ depending on platform.
Measles reaches 0.5 at $k=3$ on expedition ships and $k=5$ on the other classes, with the 0.9
threshold occurring only at higher $k$ for several large-ship curves.
SARS-CoV-2 is flatter on the large-ship surfaces in this response setting; among the supplied curves,
only the expedition surface crosses 0.5.

\section{Cross-pathogen value of information}
Figure~\ref{fig:ch7-cross-pathogen-voi} shows the corrected cross-pathogen decision surfaces.
The most consistent result is the platform gradient.
Expedition ships have the highest value of information, and the value does not become negative over
the tested prevalence grid for P4.
The peak P4 value is \$4,914 per passenger for norovirus, \$6,072 for influenza, \$10,411 for
SARS-CoV-2, and \$10,066 for measles on expedition ships.
These peaks occur at higher prevalence on expedition ships than on larger platforms because a small
absolute number of introductions can move the whole voyage across a high-cost response threshold.

\begin{figure}[htbp]
\centering
\includegraphics[width=0.92\textwidth]{figures/fig_ch7_definitive_cross_pathogen.png}
\caption{Corrected cross-pathogen value-of-information analysis for pre-boarding screening. Values compare screening policies against P0 no screening under the mid-cost scenario. P4, incentivized wearable adoption plus confirmatory testing, has the highest peak value and is the preferred active policy across most decision cells; large ships show rollover when outbreak probability saturates.}
\label{fig:ch7-cross-pathogen-voi}
\end{figure}

SARS-CoV-2 has the highest peak screening value despite lower VSP-trigger probability on large
platforms.
The reason is that the expedition curve remains highly sensitive to introductions across the tested
prevalence range, while the cost bundle attached to a triggered voyage is large.
Influenza has the narrowest screening window on large ships.
Its outbreak surfaces saturate at very low $k$, so screening is valuable at low prevalence but loses
value rapidly once multiple infectious people are likely to board.
Measles has a wider screening window than influenza because the modeled population has high
non-susceptibility, here 92\%, so introductions do not saturate the response surface as quickly.
Norovirus is intermediate: it produces positive P4 values across the full expedition grid, but the
large-ship positive window closes by approximately 0.27--0.57\% prevalence depending on class.

\begin{table}[htbp]
\centering
\caption{P4 value-of-information peaks under the mid-cost scenario. Prevalence is the pre-boarding infectious prevalence at the peak. Values are dollars per passenger relative to P0 no screening.}
\label{tab:ch7-p4-voi}
\scriptsize
\begin{tabular}{lrrrrrrrr}
\toprule
& \multicolumn{2}{c}{Expedition} & \multicolumn{2}{c}{Classic} & \multicolumn{2}{c}{Spirit} & \multicolumn{2}{c}{Mega} \\
\cmidrule(lr){2-3}\cmidrule(lr){4-5}\cmidrule(lr){6-7}\cmidrule(lr){8-9}
Pathogen & Peak VoI & Prev. & Peak VoI & Prev. & Peak VoI & Prev. & Peak VoI & Prev. \\
\midrule
Norovirus & 4914 & 1.207\% & 512 & 0.125\% & 298 & 0.097\% & 185 & 0.045\% \\
SARS-CoV-2 & 10411 & 2.000\% & 284 & 0.342\% & 232 & 0.206\% & 148 & 0.125\% \\
Influenza & 6072 & 0.567\% & 679 & 0.058\% & 469 & 0.045\% & 297 & 0.027\% \\
Measles & 10066 & 2.000\% & 729 & 0.440\% & 479 & 0.266\% & 330 & 0.160\% \\
\bottomrule
\end{tabular}
\end{table}

P4 is not uniformly best in every row of the corrected grid.
It is the top active policy in 258 of 320 pathogen--platform--prevalence cells.
The remaining 62 cells are mostly high-prevalence large-ship endpoints where all active policies have
negative value and the less intensive P2 policy loses less money.
This is not a recommendation to use P2 at high prevalence.
It means that once the rollover region is reached, the better decision may be to forgo screening or
change the voyage-level action rather than operate a screening programme with poor marginal return.

\section{Policy recommendations by platform class}
Expedition ships should treat pre-boarding screening as a default risk-management tool for all four
pathogens evaluated here.
Their P4 value remains positive across the tested prevalence grid, and the maximum value is measured
in thousands of dollars per passenger.
Small absolute changes in $K_{\mathrm{board}}$ have large consequences because the voyage has less
capacity to absorb introductions before crossing a response threshold.

Classic and spirit-class ships should use P4 selectively when pre-boarding prevalence is in the
positive window.
For norovirus, the positive P4 window extends to about 0.57\% on classic ships and 0.44\% on spirit
ships.
For influenza, the window is narrower, ending near 0.44\% on classic ships and 0.27\% on spirit
ships.
For measles and SARS-CoV-2, the window is wider than influenza because the surfaces roll over later.

Mega ships require the most conservative screening trigger.
The peak P4 value is positive for all four pathogens, but the positive window closes earliest:
approximately 0.27\% for norovirus, 0.34\% for SARS-CoV-2, 0.16\% for influenza, and 0.57\% for
measles.
At higher prevalence, the relevant operational choice may shift from individual screening to voyage
modification, targeted cancellation, pre-boarding delay, or a pre-departure testing mandate with a
different cost structure.

\section{Sensitivity to wearable accuracy and cost assumptions}
The corrected Monte Carlo outputs use the mid-cost scenario and pathogen-specific wearable accuracy.
They do not include a stored low-cost, high-cost, or full sensitivity grid.
The model code defines those parameters, so the direction of sensitivity is transparent even where
additional runs are still needed.
Higher wearable sensitivity reduces $K_{\mathrm{board}}$ and should increase value when the outbreak
surface is steep.
Higher specificity reduces false-positive denied boarding and is most important for low-prevalence
voyages, where most flagged people are not infectious.
This specificity dependence is visible in the policy comparison: P3, which denies boarding on the
wearable flag alone, is penalized by false positives, while P2 and P4 use confirmatory testing to
control that cost.

Cost sensitivity is also asymmetric.
Raising VSP, outbreak, and reputational costs increases the value of preventing a trigger.
Raising screening, secondary-testing, and denied-boarding costs narrows the positive prevalence
window.
The current cost scenario sets the VSP-related bundle at \$4.5 million per triggered voyage when
outbreak, VSP, and reputation terms are combined.
Under that assumption, boundary screening is most attractive where one or two intercepted
introductions can change the trigger probability.
Future sensitivity runs should report the break-even prevalence under the low, mid, and high cost
scenarios already present in the model configuration.

\section{Implication for intervention design}
The intervention programme should be split into two layers.
Mid-voyage surveillance should be retained as a response-quality and situational-awareness layer,
with expectations calibrated to the saturation result in Table~\ref{tab:ch7-intervention-comparison}.
Pre-boarding screening should be treated as the higher-value decision layer because it moves the
voyage along the $K_{\mathrm{board}}$ outbreak surface before amplification begins.
For the tested parameterization, P4 is the preferred operational starting point: broad incentivized
adoption, confirmatory testing before denied boarding, and platform-specific prevalence thresholds
that prevent screening from operating after value has rolled over.
